\section{Limitations and threats to validity}

Ratings and recommendation flags are noisy proxies for sentiment. The benchmark covers three Korean review domains and binary labels, so the findings may not transfer to other genres, label schemes, or languages. NSMC uses an official test file while Shopping and Steam use seeded random splits.

All neural representations are learned from scratch. This isolates data scaling but does not answer how a pretrained Korean language model would behave. Each dataset-model pair receives the same number of tuning candidates, not the same compute budget, and selection occurs at 8,000 labels. A configuration that is optimal at that size may not be optimal at 200 labels or full data.

Three final seeds describe variation but provide little power for formal tests. We therefore report effect sizes, sample SDs, and raw paired seed differences without significance language. CUDA training and CPU classical timing reflect a practical mixed-device setup. They do not isolate architecture from hardware. PyTorch deterministic mode emitted a warning for one attention kernel, so the process is seeded and repeatable at the protocol level but is not claimed to be bitwise deterministic.

The bonus adds 4,000 source examples, uses one Transformer backbone, and evaluates only four target sizes. It tests whether existing source labels help, not whether source labeling is cost-free or whether total annotation decreases. Projection diagnostics use fixed deterministic samples and cosine pairs, but stored embeddings were not retained for alternative projections.
